\documentclass[12pt, a4paper]{article}
\usepackage{bookmark}
\usepackage{amsmath}
\usepackage{amsthm}
\usepackage{amssymb}
\usepackage{tikz}
\usepackage{pgfplots}
\usepackage[utf8]{inputenc}
\usepackage{amsfonts}
\usepackage{nicematrix}
\usepackage[x11names]{xcolor}
\usepackage{geometry}
\usepackage{graphicx}
\usepackage{graphics}
\usepackage{indentfirst}
\usepackage[export]{adjustbox}
\usepackage{fancyhdr}
\usepackage[portuguese]{babel}
\usepackage{hyperref}
\usepackage{multirow}
\usepackage{setspace}
\usepackage{lastpage}
\usepackage{mathtools}
\usepackage[many]{tcolorbox}
\usepackage{newtxsf}
\usepackage{subfiles}
\usepackage{flafter}
\usepackage{float}
\usepackage{accents}
\usepackage[T1]{fontenc} 

 \setlength{\headheight}{14.5pt}
 \setlength{\parindent}{1.25cm}
 \setlength{\parskip}{0pt}
 \onehalfspacing

 \pagestyle{fancy}

 \pgfplotsset{compat = 1.18}

 \hypersetup{
     colorlinks,
     citecolor=black,
     filecolor=black,
     linkcolor=black,
     urlcolor=black
 }
 \newtheorem*{def*}{\underline{Definição}}
 \newtheorem*{theorem*}{\underline{Teorema}}
 \newtheorem*{lemma*}{\underline{Lema}}
 \newtheorem*{prop*}{\underline{Proposição}}
 \newtheorem{example}{\underline{Exemplo}}
 \theoremstyle{definition}
 \newtheorem*{proof*}{\underline{Prova}}
 \newtheorem*{crl*}{\underline{Corolário}}
 \newtheorem{exr}{\underline{Exercício}}
 \renewcommand\qedsymbol{$\blacksquare$}

\newcommand{\Hc}{\mathcal{H}}
\newcommand{\R}{\mathbb{R}}
\newcommand{\N}{\mathbb{N}}
\newcommand{\eps}{\varepsilon}


% \rfoot{Página \thepage \hspace{1pt} de \pageref{LastPage}}

 \geometry{a4paper, left=3cm, top=3cm, right=2cm, bottom=2cm}

\begin{document}
\begin{figure}[ht]
	\minipage{0.76\textwidth}
	\includegraphics[width=4cm]{../../../icmc.png}
	\hspace{7cm}
	\includegraphics[height=4.9cm,width=4cm]{../../../brasao_usp_cor.jpg}
	\endminipage
\end{figure}

\begin{center}
	\vspace*{1.0cm}
	{\Large UNIVERSIDADE DE SÃO PAULO}\\[0.4cm]
	{\Large INSTITUTO DE CIÊNCIAS MATEMÁTICAS E COMPUTACIONAIS -- ICMC}\\[1.8cm]

	{\Large\bfseries Espaços Métricos, Sistemas Iterados de Funções e Recorrência}\\[0.35cm]

	{\normalsize Relatório de atividades}\\[1.4cm]

	\begin{flushleft}
		\textbf{Autor:} Renan Wenzel\\[0.25cm]
		\textbf{Orientador:} Tiago Pereira da Silva\\[0.25cm]
		\textbf{Instituição:} Instituto de Ciências Matemáticas e Computacionais -- USP\\[0.25cm]
		\textbf{Tema de estudo:} fractais, sistemas iterados de funções, recorrência e teoria ergódica
	\end{flushleft}

	\vfill
	São Carlos, 2026
\end{center}

\newpage

\section*{Revisando Espaço de Compactos}
Durante este semestre, houve a inclusão de um novo integrante, Igor, que começou a estudar as bases para entender sistemas iterados de funções, jogos do caos e teoria ergódica.

Tendo em vista que ele não possuía uma forte base nos métodos da matemática e de conceitos importantes, foi julgado necessário a retomada do assunto do início em seminários apresentados especialmente por ele, com pontuações, dicas e correções onde julgou-se necessário, e por conta disso, a maior parte dos assuntos tratados no período do relatório parcial repetem os temas previamente vistos.
Após finalizada a parte introdutória, o grupo começou o estudo da parte de teoria ergódica, utilizando o conjunto de Cantor como um exemplo.

Recapitulando brevemente o que foi estudado, o estudo começou definindo o espaço de conjuntos compactos não vazios de um espaço métrico \((X, d)\), denotado por \(\mathcal{H(X)}\), junto à formalização da noção de distância de um ponto a um conjunto, que levou à primeira ideia de uma distância entre conjuntos, mas que falhava em ser simétrica.
Em seguida, a distância de Hausdorff, métrica oficial entre compactos foi construída, analisada e determinada sobre suas principais propriedades, culminando no espaço principal onde estarão os conjuntos autossimilares: \((\mathcal{H}(X), h)\), onde h é a métrica
\[
	h(A, B)=\max_{}(d(A, B), d(B, A)), \quad d(A, B) = \min\limits_{x\in A}d(x, B),
\]
onde A e B são compactos não vazios. Com ela, foi confirmado que \((\mathcal{H}(X), h)\) é um espaço métrico que herda a compacidade de X.

Com o espaço básico construído, foram vistos os conceitos de dilatação dos conjuntos e como eles relacionam-se com a distância de Hausdorff; o objetivo deste assunto foi encontrar uma maneira de melhor operacionalizar a métrica de Hausdorff, permitindo a demonstração do Lema da Extensão:
\begin{lemma*}
	Seja X um espaço métrico (não necessariamente completo), \((A_{n})_{n}\) uma sequência de Cauchy em \(\mathcal{H}(X)\) e \((x_{n_{i}})_{i}\) uma sequência de elementos \(x_{n_{i}}\) em \(A_{n_{i}}\); então, existe uma sequência de Cauchy \((\tilde{x}_{n})_{n}\) de elementos de A, tais que
	\[
		\tilde{x}_{n} = x_{n}.
	\]
\end{lemma*}
Explorando a relação entre as sequências de elementos e sequências de conjuntos, foi demonstrado que \((\mathcal{H}(X), h)\) herda, também, a completude do espaço métrico X.

A partir do momento que o espaço e algumas de suas propriedades topológicas foram decididas e comprovadas, o foco tornou-se às contrações, tópico no qual foram demonstradas a existência e unicidade de pontos-fixos para contrações por meio do Teorema do Ponto-Fixo de Banach, que fora enunciado como
\begin{theorem*}[Teorema da Contração de Banach]
	Seja \(f:X\rightarrow X\) uma contração em um espaço métrico completo X; então, existe um único \(x_{f}\in X\) tal que
	\[
		f(x_{f})=x_{f},
	\]
	e, para todo \(x\in X\), a sequência de composições \(\{f^{n}(x)\}_{n}\) converge para \(x_{f},\) onde
	\[
		f^{n}(x)= \underbrace{f\circ f\circ\cdots\circ f}_{n-\text{vezes}}.
	\]
\end{theorem*}
A prova seguiu, como passos, demonstrar que a sequência é de Cauchy; que ela tem um ponto de acumulação em X; que \(x_{f}\) era de fato um ponto fixo; e a unicidade.
Alguns exemplos de aplicações foram estudados, como o caso das transformações afins e suas relações às contrações, e a demonstração de que contrações de fato tornam o conjunto \(f(X)\) num subconjunto próprio de X.

\section*{Sistemas Iterados de Funções}
Quanto à introdução de dinâmica e sistemas iterados de funções, a definição utilizada foi de
\begin{def*}
	Um \textbf{sistema iterado de funções hiperbólico}, ou IFS, é um par
	\[
		\{X;w_1,\ldots,w_N\},
	\]
	em que $X$ é um espaço métrico completo e cada
	\[
		w_i:X\to X
	\]
	é uma contração com fator $s_i$. O valor \(s=\max\limits_{1\leq i\leq N}(s_{i})\) é chamado \textbf{fator de contração do IFS} \(\square\)
\end{def*}
Para que sequer faça sentido chamar s de ``fator de contração do IFS'', que é essencialmente só um par de objetos, utiliza-se o operador associado
\[
	W:\Hc(X)\to\Hc(X),
	\qquad
	W(A)=\bigcup_{i=1}^{N}w_i(A),
\]
e usando a propriedade da distância de Hausdorff para uniões, obtém-se
\[
	h(W(A),W(B))
	\leq
	s\,h(A,B),
	\qquad
	s=\max_{1\leq i\leq N}s_i<1.
\]
Portanto, $W$ é uma contração em $\Hc(X)$.
Como $X$ é completo, $\Hc(X)$ também é completo, garantindo, pelo Teorema do Ponto Fixo de Banach aplicado diretamente a $W$, a existência de um ponto-fixo \(A\in \mathcal{H(X)}\) tal que
\[
	A = W(A) = W^{2}(A) = \dotsc ,
\]
dando sentido ao nome de um conjunto autossimilar. Além disso, utilizando a propriedade de qualquer ponto ser ``puxado'' para o ponto-fixo proveniente do Teorema da Contração de Banach, nota-se que, na distância de Hausdorff, qualquer conjunto B em \(\mathcal{H(X)}\) que sofra aplicações consecutivas de W convergirá sem falha ao ponto-fixo \(A\):
\begin{def*}
	O ponto fixo $A$ do operador $W$ é chamado \textbf{atrator do IFS}.\(\square\)
\end{def*}
A ideia de ser um atrator é justamente por conta da última constatação antes da definição, e fornece uma formalização da autossimilaridade que os conjuntos fractais apresentam. Um dos exemplos mais relevantes é o conjunto de Cantor:
\begin{example}[Conjunto de Cantor]
	A construção usual do conjunto de Cantor ilustra como ele pode ser um candidato a atrator: considere
	\[
		w_1(x)=\frac{x}{3},
		\qquad
		w_2(x)=\frac{x}{3}+\frac{2}{3};
	\]
	para $A_0=[0,1]$,
	\[
		\begin{aligned}
			W(A_0)
			 & =w_1([0,1])\cup w_2([0,1]) \\
			 & =\left[0,\frac13\right]
			\cup
			\left[\frac23,1\right].
		\end{aligned}
	\]

	Na segunda aplicação de $W$, cada intervalo é novamente substituído por duas cópias contraídas ao terço, tal que, repetindo, obtém-se uma sequência de compactos
	\[
		A_0,\ W(A_0),\ W^2(A_0),\ldots
	\]
	que converge, na distância de Hausdorff, para o atrator
	\[
		C=\lim_{n\to\infty}W^n([0,1]),
	\]
	isto é, o conjunto de Cantor, justamente porque é a consequência de infinitas pulverizações em terços.
	Esse exemplo também é interessante por tornar explícita a relação entre a propriedade de autossimilaridade e a equação de ponto fixo
	\[
		C=w_1(C)\cup w_2(C).
	\]
\end{example}

A última atividade realizada teve como objetivo dar motivação ao estudo futuro de teoria ergódica, abordando os aspectos das transformações preservando medidas, ergodicidade e médias temporais, e um resultado inicial

Para relacionar essas ideias ao conjunto de Cantor, considere inicialmente a transformação mapeando x para 3x:
\[
	T:[0,1)\longrightarrow[0,1),
	\qquad
	T(x)=3x-\lfloor 3x\rfloor.
\]
Se
\[
	x=0.a_1a_2a_3\ldots{}_3 == \frac{a_1}{3} + \frac{a_2}{3^{2}}+ \frac{a_3}{3^{3}} + \dotsc
\]
é a expansão ternária de $x$, então, usando a expansão que não termina em uma sequência infinita de 2
\begin{align*}
	T(x) = 3x - a_1 & = 3\biggl(\frac{a_1}{3} + \frac{a_2}{3^{3}}+\dotsc \biggr)-a_1 \\
	                & = a_1 + \frac{a_2}{3} + \frac{a_3}{3^{2}}+\dotsc -a_1          \\
	                & = \frac{a_2}{3} + \frac{a_3}{3^{2}} + \dotsc                   \\
	                & = 0.a_2a_3a_4\ldots{}_3
\end{align*}
Assim, $T$ atua como um deslocamento dos dígitos da expansão ternária.

O conjunto de Cantor do terço médio pode ser caracterizado como os números de \([0, 1]\) que possuem uma expansão em base 3 sem o dígito 1:
\[
	C=\left\{0.a_1a_2a_3\ldots{}_3:\;a_i\in\{0,2\}\right\},
\]
que mostra a compatibilidade entre a dinâmica de $T$ e o conjunto de Cantor: ao remover o primeiro dígito de uma sequência formada apenas por $0$'s e $2$'s, o resultado será novamente uma sequência formada apenas por $0$ e $2$.

Além disso, existe um ponto cuja órbita se acumula em todo o conjunto de Cantor, funcionando como um atrator; para construí-lo, considere o número cuja expansão ternária é obtida concatenando, em ordem crescente de tamanho, todas as combinações tipo palavras, mas formadas pelos dígitos $0$ e $2$:
\[
	x=0.\,0\,2\,00\,02\,20\,22\,000\,002\,020\,022\,200\,202\,220\,222\ldots{}_3.
\]
Se $y=0.b_1b_2b_3\ldots{}_3\in C$, então, para cada $m$, a palavra $b_1b_2\ldots b_m$ aparece na expansão de $x$, de modo que, se k é a posição onde ela aparece,
\[
	T^{k}(x) = 0.b_1b_2\ldots b_mc_{m+1}c_{m+2}\dotsc_{3}
\]
Observando que, como os dígitos no conjunto de Cantor são 0 ou 2, vale a desigualdade \(| c_{j}-b_{j} |\leq 2\), tem-se
\begin{align*}
	\left|T^{k}(x)-y\right| = \biggl\vert \sum\limits_{j=m+1}^{\infty}\frac{c_{j}-b_{j}}{3^{j}} \biggr\vert & \leq \sum\limits_{j=m+1}^{\infty}\frac{2}{3^{j}}            \\
	                                                                                                        & =2\biggl(\frac{1}{3^{m+1}}+\frac{1}{3^{m+2}}+\dotsc \biggr) \\
	                                                                                                        & \leq \frac{1}{3^m},
\end{align*}
e, portanto, $T^{k}(x)\to y$. Logo, a órbita de $x$ é densa em $C$ e, em particular, acumula-se no conjunto de Cantor.

Após as bases serem estabelecidas, os estudos voltaram-se à introdução da teoria ergódica, que generaliza e sistematiza a construção apresentada para a origem do conjunto de Cantor como uma órbita se acumulando; pouco antes do teorema do ponto fixo de Banach, um seminário de motivação fora dado, onde apresentou-se a ideia de pontos recorrentes e medidas invariantes
\begin{def*}
	Se $(M,\mathcal{B},\mu)$ é um espaço de medida e $f:M\to M$ é mensurável, diz-se que $\mu$ é $f$-invariante quando, para todo $E\in\mathcal{B}$,
	\[
		\mu(E)=\mu(f^{-1}(E)). \square
	\]
\end{def*}
Com a invariância, as ideias de permanência e recorrência foram introduzidas, respectivamente sob a forma de um teorema e a definição
\begin{theorem*}
	Se $f:M\to M$ é mensurável, $\mu$ é uma medida finita e invariante, e $E\subseteq M$ é mensurável com $\mu(E)>0$, então, para $\mu$-quase todo $x\in E$, existem infinitos $n\in\N$ tais que
	\[
		f^n(x)\in E.
	\]
\end{theorem*}
A ideia da prova consiste em considerar o conjunto de pontos que saem de $E$ e nunca mais retornam; as pré-imagens iteradas desse conjunto são disjuntas e possuem, por invariância, a mesma medida. Sendo assim, como a medida total do espaço é finita, esse conjunto deve ter medida nula. A finalização consiste em mostrar que o conjunto de pontos que retornam apenas finitas vezes também possui medida nula.
\begin{def*}
	Um ponto $x$ de um espaço topológico $M$ é \textbf{recorrente para $\mathbf{f}$} quando existe uma sequência de inteiros positivos $n_1,n_2,\ldots$ tal que
	\[
		f^{n_i}(x)\longrightarrow x. \square
	\]
\end{def*}
Quando o espaço possui uma base enumerável e a medida é finita e invariante, todo ponto será recorrente.

Em estudos futuros, será construída uma medida de probabilidade invariante suportada no conjunto de Cantor $C$. De forma breve, pode-se atribuir probabilidade $1/2$ a cada escolha dos dígitos $0$ e $2$, essencialmente oferecendo 50\% de chance de aleatoriamente escolher qualquer um deles e, para cada intervalo dado por uma palavra $a_1\ldots a_n\in\{0 ,2\}^n$, definir a medida
\[
	\mu(I_{a_1\ldots a_n})=2^{-n}.
\]
Como resultado, é obtida uma medida de probabilidade de Borel $\mu$ suportada em $C$, e como a transformação $T$ remove o primeiro dígito da expansão ternária e os dígitos são escolhidos com a mesma distribuição em cada posição, essa medida satisfaz
\[
	\mu(T^{-1}(E))=\mu(E),
\]
para todo conjunto de Borel $E$ no intervalo \([0, 1]\); portanto, $\mu$ é $T$-invariante.
\begin{def*}
	Uma transformação que preserva medida $(X,\mathcal{B},\mu,T)$ é \textbf{ergódica} quando todo conjunto invariante $E$ satisfaz
	\[
		\mu(E)=0
		\qquad\text{ou}\qquad
		\mu(X\setminus E)=0. \square
	\]
\end{def*}

\subsection*{Conclusões e Direções Futuras}
Com os temas estudados, tanto o grupo quanto o integrante novo ficaram mais próximos de um nível semelhante para dar continuidade aos estudos na área de teoria ergódica, além da motivação fornecida pelos exemplos mais avançados que guiam estudos futuros.

Dentre os possíveis temas de futuros estudos, estão os \textit{Martingales}, a construção de uma medida suportada no conjunto de Cantor, e explorar a teoria ergódica mais profundamente.

\begin{thebibliography}{9}
	\bibitem{barnsley}BARNSLEY, Michael F. Fractals everywhere. 2. ed. \textbf{Boston: Academic Press}, 1993.
	\bibitem{tao}TAO, Terence. An introduction to measure theory. \textbf{Providence: American Mathematical Society}, 2011. (Graduate Studies in Mathematics, v. 126).
\end{thebibliography}

\end{document}
